% Teil II: durchgerechneter Weg, Teile a--c.

\subsection*{Gegeben und gesucht}
Ballon: Hülle $m_{\mathrm{H}}=\qty{10}{kg}$, $V_{\max}=\qty{1000}{m^3}$.
$R=\qty{8.314}{J/(mol\,K)}$, $T_0=\qty{273.15}{K}$,
$p_0=\qty{1.013e5}{Pa}$; $g=\qty{9.81}{m/s^2}$; molare Massen
$M_{\mathrm{He}}=\qty{4.0}{g/mol}$, $M_{\mathrm{N}_2}=\qty{28.0}{g/mol}$,
$M_{\mathrm{O}_2}=\qty{32.0}{g/mol}$, $M_{\mathrm{Ar}}=\qty{39.9}{g/mol}$;
Luft $78\%$ N$_2$, $21\%$ O$_2$, $1\%$ Ar.
Gesucht: (a) Dichten von He und Luft; (b) Mindest-Heliumvolumen zum Abheben;
(c) H"ohe, in der $V_{\max}$ erreicht wird, und ob der Ballon weiter steigt.

\subsection*{Ansatz und Physik dahinter}
Dichte aus idealem Gas $\rho=Mp/(RT)$. Effektive molare Masse der Luft aus
den gewichteten Anteilen. Auftriebsgleichgewicht am Boden: verdr"angte
Luftmasse $=$ Hülle $+\,$ Helium; daraus $V_0$. Steigt der Ballon, bleibt die
Heliummenge konstant und $pV=\text{const}$ (isotherm), also
$V(h)=V_0\,p_0/p(h)$; $V=V_{\max}$ liefert $h$.

\subsection*{Rechnung}
\textbf{a) Dichten.}
\begin{equation}
  M_{\mathrm{He}}=4.0\times10^{-3}\ \mathrm{kg/mol},\qquad
  \rho_{\mathrm{He}}=\frac{p_0 M_{\mathrm{He}}}{R T_0}
  =\frac{1.013\times10^{5}\cdot4.0\times10^{-3}}{8.314\cdot273.15}
  \approx\boxed{\;\qty{0.18}{kg/m^3}\;}.
\end{equation}
Effektive Luftmolmasse:
\begin{equation}
  M_{\mathrm{Luft}}=0.78\cdot28+0.21\cdot32+0.01\cdot39.9
  =28.96\ \mathrm{g/mol}\approx\qty{0.02896}{kg/mol}\;,
  \qquad
  \rho_{\mathrm{Luft}}=\frac{p_0 M_{\mathrm{Luft}}}{R T_0}
  \approx\boxed{\;\qty{1.29}{kg/m^3}\;}.
\end{equation}

\textbf{b) Mindestvolumen zum Abheben.} Auftrieb $=$ Gewicht:
\begin{equation}
  \rho_{\mathrm{Luft}}\,V_0\,g=(m_{\mathrm{H}}+\rho_{\mathrm{He}}V_0)\,g
  \quad\Rightarrow\quad
  V_0=\frac{m_{\mathrm{H}}}{\rho_{\mathrm{Luft}}-\rho_{\mathrm{He}}}
  =\frac{10}{1.29-0.18}\approx\boxed{\;\qty{9.0}{m^3}\;}
  \ (\approx \qty{8.98}{m^3}).
\end{equation}
Gegenprobe: Hülle 10~kg, Helium 0.18$\cdot$9.0$\approx$1.6~kg, zusammen
$\approx$11.6~kg; verdr"angt 1.29$\cdot$9.0$\approx$11.6~kg Luft. Gleich-
gewicht. \checkmark

\textbf{c) H"ohe mit $V=V_{\max}$; steigt er weiter?} Bei koh"arenter
isothermer F"ullung $V(h)=V_0\,p_0/p(h)$:
\begin{equation}
  V_{\max}=V_0\,e^{+M_{\mathrm{Luft}}\,g\,h/(R T_0)}
  \;\Rightarrow\;
  h=\frac{R T_0}{M_{\mathrm{Luft}}\,g}\ln\frac{V_{\max}}{V_0}.
\end{equation}
Mit $M_{\mathrm{Luft}}=0.02896$~kg/mol:
\begin{equation}
  \frac{R T_0}{M_{\mathrm{Luft}}\,g}
  =\frac{8.314\cdot273.15}{0.02896\cdot9.81}
  \approx\qty{7991}{m},\qquad
  h=\qty{7991}{m}\cdot\ln(1000/8.98)=\qty{7991}{m}\cdot4.712
  \approx\boxed{\;\qty{37.6}{km}\;}.
\end{equation}
\textbf{Weiter steigen?} In Schritt (c) ist das Volumen am Maximum: Der
Ballon k"onnte an dieser H"ohe nicht weiter aufbl\"ahen ($V_{\max}$ ist das
Fassungsvermögen der Hülle, das Helium hat den Umgebungsdruck erreicht).
Oberhalb davon f"allt $p$ weiter, der Ballon k"onnte nicht weiter expandieren,
zugleich nimmt die verdr\"angte Luftmasse $\rho_{\mathrm{Luft}}V_{\max}$ mit
der H"ohe ab, w"ahrend das Gewicht (Hülle + Helium) konstant bleibt:
der Auftriebsspielraum schrumpft, das Gleichgewicht wird bei $V_{\max}$
stabil erreicht. Physikalisch hat er seine ``Reiseh"ohe'' erreicht, wo die
verdr\"angte Luftmasse wieder dem Gewicht entspricht --- ein stabiler
Schwebezustand. Er steigt also \emph{nicht} weiter.

\begin{ergebnis}
a) $\rho_{\mathrm{He}}=\qty{0.18}{kg/m^3}$, $\rho_{\mathrm{Luft}}=\qty{1.29}{kg/m^3}$,
$M_{\mathrm{Luft}}=\qty{29.0}{g/mol}$.
b) $V_0=m_{\mathrm{H}}/(\rho_{\mathrm{Luft}}-\rho_{\mathrm{He}})\approx\qty{9.0}{m^3}$.
c) $h=\dfrac{RT_0}{M_{\mathrm{Luft}}g}\ln\dfrac{V_{\max}}{V_0}
\approx\qty{37.6}{km}$; der Ballon steigt nach Erreichen von $V_{\max}$
nicht weiter (stabiler Schwebezustand).
\end{ergebnis}

\begin{kontrolle}
\textbf{Einheiten:} $\rho\colon\dfrac{Pa\cdot kg/mol}{J/(mol\,K)\cdot K}
=\dfrac{kg/m^2}{m/s^2\cdot s^{-1}}$... konkret:
$pM/(RT)=(\mathrm{kg\,m/s^2})/(\mathrm{m^2})\cdot(\mathrm{kg/mol})
/[(\mathrm{kg\,m^2/s^2})/\mathrm{mol}]=\mathrm{kg/m^3}$. \checkmark
$h$: $RT_0/(Mg)=[\mathrm{m^2/s^2}]/[\mathrm{m/s^2}]=\mathrm m$. \checkmark

\textbf{Größenordnung/Grenzfälle:} Helium 0.18 vs.\ Luft 1.29 kg/m$^3$ ---
Helium ist ~7"mal leichter; $V_0=9$~m$^3$ (weniger als 2~m Kugeldurchmesser)
trägt die 10-kg-Hülle knapp. 37.6~km ist die typische H"ohe von
Wetterballons (30--40~km, Stratosph"are) --- die Antwort ist physikalisch
plausibel. Bei $T_0\to0$ ginge $\rho\to\infty$, $V_0\to0$: unrealistisch,
spiegelt nur das Modell. \checkmark
\end{kontrolle}

\begin{merksatz}
Ideales Gas: $\rho=Mp/(RT)$. Auftrieb $=\rho_{\mathrm{Luft}}Vg$; Abheben
genau bei $V=m_{\mathrm{H}}/(\rho_{\mathrm{Luft}}-\rho_{\mathrm{He}})$.
Isothermer Ballon: $pV$ konstant, $V(h)=V_0e^{Mgh/(RT_0)}$; $V_{\max}$
erreicht er bei $h=RT_0/(Mg)\ln(V_{\max}/V_0)$ und schwebt dort stabil.
\end{merksatz}
